% Article-class smoke-build of the Chiyoda paper.
% Use this when acmart is not installed:
%   make smoke   ->   build/main_smoke.pdf

\documentclass[10pt,a4paper,twocolumn]{article}
\usepackage[utf8]{inputenc}
\usepackage[T1]{fontenc}
\usepackage{lmodern}
\usepackage[margin=0.7in]{geometry}
\setlength{\columnsep}{0.3in}
\usepackage{xcolor}
\usepackage{hyperref}
\usepackage{booktabs}
\usepackage{graphicx}
\usepackage{amssymb}
\usepackage{amsmath}
\usepackage{microtype}
\hypersetup{colorlinks=true,linkcolor=blue!50!black,citecolor=blue!50!black,urlcolor=blue!50!black}

\newcommand{\todo}[1]{{\color{red}\textbf{[TODO]}\ \emph{#1}}}
\providecommand{\subtitle}[1]{}

\IfFileExists{stats.tex}{\input{stats}}{%
  \newcommand{\statStudyName}{TBD}%
  \newcommand{\statVariantCount}{0}%
  \newcommand{\statRunCount}{0}%
  \newcommand{\statInterventionEvents}{0}%
  \newcommand{\statBestPolicy}{TBD}%
  \newcommand{\statBestEfficiency}{TBD}%
  \newcommand{\statPeakHarmfulConvergence}{TBD}%
}

\title{Chiyoda: Entropy-Guided Information Control for Hazard-Coupled Pedestrian Evacuation\\
\large A Simulation Framework for Measuring When Emergency Communication Helps or Harms}
\author{\href{https://gabrielongzm.com}{Gabriel Ong Zhe Mian}\\\texttt{gabrielzmong@gmail.com}\\Independent, Singapore}
\date{\today}

\begin{document}
\maketitle

\begin{abstract}
Emergency evacuation is usually modeled as a physical routing problem, with
information appearing as a fixed assumption. We present \emph{Chiyoda}, an
Information-Theoretic Evacuation Dynamics framework that treats emergency
communication as a controllable intervention in a coupled physical-information
system. Agents carry probabilistic beliefs, exchange distorted information
through local gossip, receive beacon and responder messages, suffer
physiological impairment under hazards, and choose routes through
belief-weighted path planning. The framework compares static, global,
responder-relay, entropy-targeted, density-aware, exposure-aware, and
bottleneck-avoidance policies. The central claim is not that more information
is always better, but that emergency communication has measurable safety
tradeoffs.
\end{abstract}

\section{Introduction}
\label{sec:intro}

Emergency communication is usually treated as a benign input to evacuation:
give people better information and their decisions should improve. That
assumption is plausible at the individual level, but it is not obviously true
at the crowd level. A warning that reduces uncertainty for everyone at once can
also synchronize route choice, push agents toward the same exit, amplify a
bottleneck, or move people through a contaminated region before the physical
hazard has dissipated. The question is therefore not simply whether information
helps, but when, where, and for whom it helps.

This paper studies evacuation as a coupled physical-information system. The
physical side includes layout geometry, social-force movement, bottleneck
dynamics, hazard spread, and physiological impairment. The information side
includes partial exit knowledge, hazard beliefs, signage and PA broadcasts,
responder credibility, local gossip, message decay, and distortion. Chiyoda
connects these layers by making agents plan through their believed world rather
than the ground truth world. They may avoid a hazard they have heard about,
miss an exit they do not know, herd when uncertainty is high, or follow a
stale route because the belief has not yet decayed.

\paragraph{Research question.} We ask whether entropy-guided communication
can improve evacuation outcomes without creating harmful convergence. More
concretely: which intervention policies reduce belief entropy and improve
belief accuracy per unit of induced exposure, queueing, and exit imbalance?
The current study scaffold compares static signage, global broadcasts,
responder relay, entropy-targeted broadcasts, density-aware targeting,
exposure-aware targeting, and bottleneck-avoidance targeting.

\paragraph{Contributions.} This paper makes four contributions.

\begin{enumerate}
\item \textbf{A controllable information-intervention layer.} Chiyoda turns
emergency communication into a first-class policy object rather than a fixed
scenario assumption. Policies can be static, authority-driven, or adaptive to
entropy, density, exposure, and bottleneck pressure.

\item \textbf{A coupled ITED simulation loop.} Agents carry probabilistic
beliefs over exits and hazards, update them through observation, beacon
broadcast, responder relay, and gossip, then route through belief-weighted
costs under social-force dynamics and hazard impairment.

\item \textbf{Information-safety metrics.} We introduce study outputs that
jointly measure entropy reduction, belief accuracy gain, intervention reach,
hazard exposure pressure, queue pressure, exit imbalance, and a harmful
convergence index.

\item \textbf{A reproducible paper pipeline.} YAML study definitions generate
multi-seed bundles with Parquet/CSV tables and paper figures. The LaTeX paper
ingests generated study statistics through \texttt{stats.tex}.
\end{enumerate}

\paragraph{Roadmap.} \S\ref{sec:background} positions the work against
pedestrian dynamics, information-aware evacuation, and ABM validation work.
\S\ref{sec:design} describes Chiyoda's model and intervention policies.
\S\ref{sec:implementation} explains the software pipeline. \S\ref{sec:evaluation}
lays out the planned empirical study. \S\ref{sec:related} sharpens the
distinction from prior work. \S\ref{sec:limitations} states the current
threats to validity, and \S\ref{sec:conclusion} concludes.

\section{Background}
\label{sec:background}

Pedestrian evacuation models have a long physical tradition. The social-force
model represents pedestrians as self-driven particles subject to destination,
repulsion, and contact forces~\cite{helbing1995social}. Cellular automata
models discretize space and update local movement decisions through transition
rules~\cite{burstedde2001simulation}. Evacuation-model reviews emphasize that
these physical models must be interpreted together with scenario assumptions,
calibration data, and validation scope rather than as universally predictive
engines~\cite{gwynne1999review,schadschneider2009evacuation}.
Fundamental-diagram work relates crowd density to walking speed and flow
capacity~\cite{fruin1971pedestrian,weidmann1993transporttechnik}.
These tools are useful because evacuation failure is often physical:
bottlenecks, counter-flow, density waves, and slower-is-faster effects can
dominate individual intent.

Hazard-coupled evacuation adds another layer. In CBRN or smoke scenarios, the
environment changes while people move through it; visibility, speed, and
physiology can degrade as a function of exposure. Prior toxic-gas evacuation
models already show that information perception and transmission matter for
route choice under a spreading contaminant~\cite{makmul2020cellular,liu2021toxicgas}.
Chiyoda builds on this insight but makes the communication policy itself the
experimental object.

Information-aware evacuation models also intersect with route-choice theory,
game theory, and bounded rationality. Pedestrian route-choice models connect
activity constraints, perceived costs, and network structure~\cite{hoogendoorn2004routechoice},
while evacuation experiments show that visibility can change route-selection
behavior~\cite{guo2012visibility}. Bayesian formulations can model incomplete information
and route-choice equilibria under congestion~\cite{mah90dynamic,wang2023bayesian}. Fire
behavior and virtual-evacuation studies also show that exit choice is shaped
by affiliation, social influence, stress, familiarity, and other occupants'
movement~\cite{sime1983affiliative,kobes2010building,bode2013human,kinateder2014social}. Social
group ABMs review the importance of communication, affiliation, and collective
behavior, while also noting persistent gaps in behavioral realism and
validation~\cite{templeton2023social,senanayake2024agent}. Chiyoda takes a
simulation-first position: it does not claim to solve empirical validation in
one paper, but it makes belief uncertainty measurable and interventions
reproducible.

The information-theoretic lens uses Shannon entropy~\cite{shannon1948mathematical}.
For each agent, Chiyoda computes entropy over exit existence, hazard knowledge,
and general danger. Global entropy is the population mean. Entropy is not a
normative objective by itself: the goal is to test when reducing entropy
improves safety and when it merely concentrates the crowd's behavior.

\section{Design}
\label{sec:design}

Chiyoda models evacuation as a feedback loop between physical state,
information state, and action. At each step, hazards evolve; agents observe
nearby exits and hazards; beacons, responders, gossip, and intervention
policies update beliefs; agents revise intentions; belief-weighted navigation
selects paths; social-force dynamics move agents; and telemetry records the
macroscopic outcome.

\subsection{Belief State}
\label{subsec:beliefs}

Each agent carries a belief vector over exits, hazards, and general danger.
Exit beliefs store existence probability, congestion estimate, freshness,
source credibility, and gossip hop count. Hazard beliefs store position,
severity estimate, radius estimate, freshness, credibility, and hop count.
Agents do not route through the ground truth environment directly; they route
through the subset of exits and hazards represented in their beliefs.

\subsection{Communication Channels}
\label{subsec:channels}

The baseline ITED channels are direct observation, beacon broadcast, and
agent-to-agent gossip. Observation is range-limited by visibility and
physiological impairment. Beacons provide high-credibility exit knowledge
within a fixed radius. Gossip is local, probabilistic, and distorted by hop
count and panic state.

\subsection{Intervention Policies}
\label{subsec:policies}

The new contribution is an intervention policy interface. Every policy decides
whether to fire, selects one or more spatial targets, builds a message, applies
it to recipients inside a radius, and emits telemetry. Chiyoda currently
implements seven policies:

\begin{itemize}
\item \textbf{Static beacon}: periodic local messages from existing signage
or PA beacon locations.
\item \textbf{Global broadcast}: station-wide high-reach communication.
\item \textbf{Responder relay}: mobile high-credibility messages emitted by
released responders.
\item \textbf{Entropy-targeted}: targets agents with highest belief entropy.
\item \textbf{Density-aware}: targets dense local clusters.
\item \textbf{Exposure-aware}: targets agents under high current or cumulative
hazard pressure.
\item \textbf{Bottleneck-avoidance}: targets active bottleneck zones.
\end{itemize}

Figure~\ref{fig:policy-taxonomy} summarizes the policy space. Figure~\ref{fig:intervention-loop}
shows where the intervention hook sits in the simulation loop.

\begin{figure}[!t]
  \centering
  \includegraphics[width=\linewidth]{figures/policy-taxonomy.pdf}
  \caption{Information intervention policy taxonomy. Policies differ by
  source, target signal, spatial reach, and safety objective.}
  \label{fig:policy-taxonomy}
\end{figure}

\begin{figure}[!t]
  \centering
  \includegraphics[width=\linewidth]{figures/intervention-loop.pdf}
  \caption{Intervention hook in the ITED simulation loop. Policies read the
  current physical-information state, update recipient beliefs, and emit
  before/after telemetry before navigation.}
  \label{fig:intervention-loop}
\end{figure}

\subsection{Information-Safety Metrics}
\label{subsec:metrics}

The core metric is information-safety efficiency:

\[
\eta_{safe} =
\frac{\Delta H^{+} + \Delta A^{+}}
     {1 + P_{exposure} + P_{queue}},
\]

where $\Delta H^{+}$ is positive entropy reduction, $\Delta A^{+}$ is positive
belief-accuracy gain, $P_{exposure}$ is recipient-weighted hazard-load
pressure, and $P_{queue}$ is bottleneck pressure observed at intervention
time. Chiyoda also records intervention reach, exit imbalance, and a harmful
convergence index that rises when entropy reduction coincides with concentrated
exit usage, queueing, and exposure.

\section{Implementation}
\label{sec:implementation}

Chiyoda is implemented as a Python research toolkit. Scenario YAML files define
layouts, cohorts, hazards, responders, information parameters, intervention
policies, seeds, and variants. The scenario manager builds a simulation by
loading the layout, constructing agents, hazards, exits, responders,
navigation, spatial indexing, behavior state machines, and optional
intervention policies.

The core implementation lives in \texttt{chiyoda/core/simulation.py}. The
simulation loop records step telemetry, agent telemetry, cell-level maps,
exit flow, bottleneck dynamics, hazard state, gossip events, and intervention
events. Study execution lives in \texttt{chiyoda/studies/runner.py}; it
materializes variants, runs each seed, aggregates tables, and exports study
bundles. Analysis and figure export live in \texttt{chiyoda/analysis/}.

The information intervention layer lives in
\texttt{chiyoda/information/interventions.py}. It is deliberately small:
policies produce target points and messages, while the shared executor applies
messages to recipients and computes before/after entropy and accuracy. This
keeps policy design separate from belief mutation and telemetry.

The paper pipeline mirrors the code pipeline. Running
\texttt{python3 -m chiyoda.cli sweep scenarios/study\_information\_control.yaml}
emits tables and figures; \texttt{paper/scripts/gen\_stats.py} converts the
exported bundle into \texttt{stats.tex}; and the LaTeX build injects those
values into the abstract and evaluation.

\todo{Add an architecture figure once the first full study bundle has been generated.}

\section{Evaluation}
\label{sec:evaluation}

We evaluate Chiyoda with one primary policy comparison and two supporting
stress studies. The primary study is
\texttt{study\_information\_control.yaml}, a 50-seed comparison of eight
emergency communication regimes: no deliberate intervention, static beacon
broadcast, global PA broadcast, responder relay, entropy-targeted
communication, density-aware communication, exposure-aware communication, and
bottleneck-avoidance communication. The exported bundle contains
\statRunCount{} runs, \statVariantCount{} variant aggregates, and
\statInterventionEvents{} intervention events. Each run uses the same
station-sarin scenario, a 120-agent commuter/tourist population, asymmetric
initial information, gossip, beacon infrastructure, three delayed responders,
and a 500-step horizon.

The two supporting studies each use 30 seeds. The intervention ablation study
varies targeting budget, timing, and bottleneck awareness. The message-quality
study varies broadcast credibility, delay, and repetition. These studies are
not replacements for the main comparison; they test whether the main
information-safety interpretation is robust to natural design choices.

\subsection{Outcome Measures}
\label{subsec:outcomes}

We report three classes of outcome. Physical outcomes include evacuated agents,
mean travel time, incapacitation count, mean hazard exposure, peak bottleneck
queue, and exit imbalance. Information outcomes include mean entropy, entropy
reduction, intervention reach, entropy reduction attributable to interventions,
and belief-accuracy gain. Coupled outcomes include information-safety
efficiency (ISE) and harmful convergence index (HCI). These coupled metrics are
the main object of the paper: they ask whether a policy buys useful belief
improvement without creating exposure or queueing pressure.

Variant tables report means across seeds. Pairwise comparisons against the
no-intervention baseline use two-sided Mann--Whitney tests over run-level
summary rows, with effect sizes used as descriptive support rather than as
standalone proof of operational superiority. This choice is conservative for
the current stage: the simulator is stochastic and stylized, so the evaluation
focuses on robust tradeoff patterns rather than on declaring a universal
winner.

\begin{figure*}[!t]
  \centering
  \includegraphics[width=\textwidth]{../out/information_control_50/figures/13_information_safety_frontier.pdf}
  \caption{Information-safety frontier across communication policies in the
  50-seed main study. Static beacons produce the highest information-safety
  efficiency, while density-aware and exposure-aware policies show the largest
  harmful-convergence penalties under this scenario.}
  \label{fig:safety-frontier}
\end{figure*}

\begin{figure*}[!t]
  \centering
  \includegraphics[width=\textwidth]{../out/information_control_50/figures/12_intervention_timeline.pdf}
  \caption{Intervention reach and entropy effects over time. Broad reach and
  useful safety effect separate: global broadcast reaches many agents, but its
  information-safety efficiency remains below local static beacons, bottleneck
  avoidance, responder relay, and entropy targeting.}
  \label{fig:intervention-timeline}
\end{figure*}

\subsection{Primary Policy Comparison}
\label{subsec:aggregate-findings}

Table~\ref{tab:aggregate} summarizes the 50-seed variant aggregates. The first
result is negative but important: no communication policy dominates the
baseline across physical evacuation outcomes. The no-intervention baseline
evacuates 10.76 agents on average. Static beacons evacuate 10.90, global
broadcast evacuates 11.00, and the remaining active policies range from 9.70
to 10.86. Mann--Whitney tests against the baseline do not show significant
improvements in evacuated count or hazard exposure for any active policy at
$n=50$ seeds. This is why the paper cannot claim that more information
monotonically improves evacuation.

\begin{table*}[!t]
\centering
\caption{Variant aggregates from the 50-seed information-control study. ISE
denotes information-safety efficiency; HCI denotes harmful-convergence index.}
\label{tab:aggregate}
\begin{tabular}{lrrrrrr}
\toprule
Policy & Evac. & Travel (s) & Exposure & Queue & ISE & HCI \\
\midrule
Static beacon & 10.90 & 10.88 & 2.900 & 2.72 & 0.0136 & 7.27 \\
Bottleneck avoidance & 9.70 & 9.76 & 2.904 & 2.84 & 0.0113 & 7.91 \\
Responder relay & 10.52 & 10.69 & 2.905 & 2.76 & 0.0102 & 6.31 \\
Entropy-targeted & 9.90 & 10.31 & 2.873 & 2.86 & 0.0091 & 7.02 \\
Global broadcast & 11.00 & 11.86 & 2.866 & 2.84 & 0.0058 & 8.28 \\
Density-aware & 10.78 & 10.63 & 2.907 & 2.80 & 0.0050 & 9.41 \\
Exposure-aware & 10.86 & 11.65 & 2.916 & 2.86 & 0.0026 & 8.33 \\
No intervention & 10.76 & 11.53 & 2.919 & 2.74 & 0.0000 & 0.00 \\
\bottomrule
\end{tabular}
\end{table*}

The second result is that communication policy still matters, even when
physical outcomes remain noisy. Static beacons are the best policy by
information-safety efficiency (\statBestPolicy{}, \statBestEfficiency{}). They
combine local reach, repeated exposure, and high credibility: enough to reduce
belief uncertainty without synchronizing the full population toward the same
capacity-constrained route. Responder relay has lower efficiency than static
beacons, but also the lowest HCI among the active policies in the main study
(6.31), which supports the idea that delayed, mobile, high-credibility sources
can be conservative safety controllers.

The third result is that wide reach is not the same as useful reach. Global
broadcast slightly improves mean evacuation count relative to the baseline
(11.00 versus 10.76) and has the lowest mean hazard exposure among the main
policies (2.866). However, it ranks fifth by ISE (0.0058) and has a high HCI
(8.28). In this scenario, the global message is useful but inefficient: it
homogenizes route guidance broadly, while local and capacity-aware policies
produce more information benefit per unit of safety pressure.

The fourth result is that adaptive targeting exposes a real control tradeoff.
Entropy targeting has the largest intervention accuracy gain in the main study
(1.76), and both entropy targeting and bottleneck avoidance significantly
reduce mean travel time against no intervention
($p=0.035$ and $p=0.014$, respectively). Yet these same policies do not improve
evacuation count, and bottleneck avoidance has a near-significant reduction in
evacuated agents ($-1.06$, $p=0.061$). The practical interpretation is that
better-informed trajectories can be shorter for agents who move, while still
failing to improve aggregate completion under capacity and hazard pressure.

\subsection{Intervention Ablation}
\label{subsec:ablation-results}

Table~\ref{tab:ablation} reports the 30-seed ablation study. The strongest
finding is budget non-monotonicity. Early low-budget entropy targeting has the
highest ISE in the ablation (0.0253), more than double early high-budget
entropy targeting (0.0092). The high-budget version improves belief accuracy
more strongly (1.89 versus 1.26), but also produces larger queue pressure and
lower information-safety efficiency. This supports the paper's core mechanism:
more certainty is not automatically safer when many agents receive similar
guidance under shared bottleneck constraints.

\begin{table*}[!t]
\centering
\caption{Supporting intervention ablation over 30 seeds.}
\label{tab:ablation}
\begin{tabular}{lrrrrrr}
\toprule
Condition & Evac. & Travel (s) & Exposure & Queue & ISE & HCI \\
\midrule
Entropy early, low budget & 10.30 & 11.74 & 2.809 & 2.93 & 0.0253 & 8.57 \\
Bottleneck late & 11.17 & 11.57 & 2.873 & 2.77 & 0.0143 & 7.69 \\
Bottleneck early & 10.40 & 10.85 & 2.778 & 2.97 & 0.0140 & 7.98 \\
Entropy early, high budget & 10.70 & 11.65 & 2.879 & 3.00 & 0.0092 & 7.08 \\
Entropy late, high budget & 11.47 & 11.80 & 2.866 & 2.80 & 0.0070 & 7.43 \\
No intervention & 11.20 & 11.84 & 2.869 & 2.67 & 0.0000 & 0.00 \\
\bottomrule
\end{tabular}
\end{table*}

The timing ablation also matters. Late entropy targeting recovers evacuation
count relative to early entropy targeting (11.47 versus 10.70 under the
high-budget setting), but its ISE is lower. Bottleneck targeting is less
sensitive to late deployment: the late bottleneck condition has slightly higher
evacuation count and slightly higher ISE than early bottleneck targeting. This
suggests that bottleneck-aware control is a better candidate for delayed
operational response, while entropy targeting needs stricter budget control.

\subsection{Message Quality}
\label{subsec:message-quality-results}

Table~\ref{tab:message-quality} reports the 30-seed message-quality study.
Credibility has the expected direction for information efficiency:
high-credibility global broadcast has the highest ISE among the
broadcast-quality conditions (0.0066), delayed high-credibility broadcast is
close behind (0.0063), medium credibility falls to 0.0034, and low credibility
falls to 0.0026. Repeating low-credibility messages frequently does not repair
this loss; it has the lowest ISE in the study (0.0009) while preserving high
harmful convergence (8.92).

\begin{table*}[!t]
\centering
\caption{Supporting message-quality study over 30 seeds.}
\label{tab:message-quality}
\begin{tabular}{lrrrrrr}
\toprule
Condition & Evac. & Travel (s) & Exposure & Queue & ISE & HCI \\
\midrule
High-credibility global & 11.13 & 12.01 & 2.889 & 2.83 & 0.0066 & 8.42 \\
Delayed high credibility & 10.80 & 10.70 & 2.856 & 2.80 & 0.0063 & 8.09 \\
Medium-credibility global & 11.53 & 10.85 & 2.862 & 2.80 & 0.0034 & 7.12 \\
Low-credibility global & 12.10 & 11.24 & 2.893 & 2.90 & 0.0026 & 9.07 \\
Frequent low credibility & 11.93 & 11.79 & 2.858 & 2.80 & 0.0009 & 8.92 \\
No intervention & 11.20 & 11.84 & 2.869 & 2.67 & 0.0000 & 0.00 \\
\bottomrule
\end{tabular}
\end{table*}

This support study strengthens the interpretation of global broadcast. Global
messages can increase completed evacuations in some parameterizations, but the
information-safety metric penalizes conditions where the same broad signal also
creates route synchronization. Message quality therefore behaves like a control
parameter, not a cosmetic parameter: credibility and timing change whether
communication improves useful beliefs or merely coordinates crowd movement.

\subsection{Hypotheses Revisited}
\label{subsec:hypotheses-results}

H1 predicted that adaptive targeting would improve information-safety
efficiency over global broadcast. The 50-seed result supports this for static
beacons, bottleneck avoidance, responder relay, and entropy targeting, all of
which exceed global broadcast by ISE. The result is weaker for density-aware
and exposure-aware targeting, which have nonzero efficiency but lower safety
efficiency than global broadcast in this scenario.

H2 predicted that global broadcast would reduce entropy quickly but risk route
synchronization. The data support the synchronization concern more strongly
than the entropy-reduction claim. Global broadcast has high reach and high HCI,
but its ISE trails more localized policies. The message-quality study further
shows that frequent low-credibility broadcast is especially inefficient.

H3 predicted that exposure-aware targeting would reduce hazard exposure.
Exposure-aware targeting does not meaningfully reduce exposure in the 50-seed
main study: its mean hazard exposure is 2.916, close to the baseline value of
2.919 and above several other policies. This hypothesis is not supported by
the present configuration.

H4 predicted that responder relay would be robust but delayed. The result is
partially supported. Responder relay is third by ISE, has the lowest HCI among
the main active policies, and avoids the extreme synchronization penalties of
some adaptive policies. Its limitation is not obvious physical superiority, but
conservative information control.

\subsection{Interpretation}
\label{subsec:interpretation}

The study supports the paper's central framing rather than a single winning
policy. Information interventions should be evaluated as safety-control
actions, not as monotonic entropy minimizers. A policy can reach many agents
and still have low safety efficiency; another can improve belief accuracy and
still fail to improve aggregate completion. The promising direction is
therefore not ``broadcast more,'' but ``broadcast credibly, locally, and with
awareness of capacity and exposure.'' Static beacon and responder relay
policies currently look like conservative baselines, while entropy targeting is
scientifically valuable because it exposes the failure mode created by
overconfident, synchronized route choice.

\subsection{Threats to Evaluation}
\label{subsec:evaluation-threats}

The current evidence is stronger than the pilot run, but still not final
external validation. The experiments use one station geometry, one hazard
scenario, one population size, and a fixed evacuation horizon. The supporting
studies vary intervention design but not architectural layout, sensor latency,
or calibrated pedestrian demographics. The HCI and ISE metrics are also model
constructs; future work should test alternative normalizations and compare
them against trajectory-level bottleneck events, drill data, or expert-coded
evacuation plans.

\section{Related Work}
\label{sec:related}

Chiyoda builds on physical pedestrian-evacuation models rather than replacing
them. Social-force, cellular-automata, and microscopic interaction models
explain how local movement rules can produce bottlenecks, counterflow, density
waves, and other crowd-level failure modes~\cite{helbing1995social,
burstedde2001simulation,moussaid2011simple}. Evacuation-model reviews caution
that such models remain engineering abstractions whose outputs depend on
scenario scope, behavioral assumptions, calibration, and validation
data~\cite{gwynne1999review,schadschneider2009evacuation}. Chiyoda uses this
tradition as the motion substrate; its distinction is the information-control
layer placed on top of that substrate.

Hazard-coupled evacuation and guidance models are closest to the CBRN-like
setting studied here. Toxic-gas and gaseous-material simulations show that
contaminant spread, perception, and information transmission can alter route
choice and evacuation outcomes~\cite{makmul2020cellular,liu2021toxicgas}.
Guidance models similarly show that recommendations must account for hazards,
congestion, and compliance~\cite{chu2017variable}. Chiyoda differs by making
message policy the controlled variable: timing, targeting, budget, credibility,
and reach are varied directly instead of being treated as fixed scenario
assumptions.

Incomplete-information and social-behavior models motivate the belief layer.
Bayesian and game-theoretic route-choice work studies uncertainty and
congestion~\cite{hoogendoorn2004routechoice,wang2023bayesian}, while
social-group ABM surveys and evacuation studies emphasize communication,
affiliation, familiar routes, visible crowd movement, and conflicting social
cues~\cite{templeton2023social,senanayake2024agent,sime1983affiliative,
kobes2010building,bode2013human,kinateder2014social}. Chiyoda keeps these
effects stylized, but exports belief entropy, belief accuracy, reach, exposure,
queueing, and exit imbalance together so belief improvement can be compared
against downstream safety effects.

Existing tools cover much of the surrounding simulation and validation stack.
JuPedSim, Vadere, and SUMO provide mature pedestrian simulation capabilities,
and PedPy provides trajectory-analysis measures for density, velocity, flow,
and fundamental diagrams~\cite{kemlohwagoum2015jupedsim,kleinmeier2019vadere,
alvarezlopez2018sumo,schrodter2025pedpy}. Geometry and hazard studies can draw
on OpenStationMap, GTFS Pathways, and high-fidelity tools such as the NIST Fire
Dynamics Simulator~\cite{openstationmap2026,gtfs2026pathways,mcgrattan2013fds}.
Chiyoda is therefore best read as a replayable information-intervention
benchmark that can be calibrated or cross-checked against these systems, not as
a substitute for full pedestrian or dispersion validation~\cite{ronchi2013verification}.

\section{Limitations}
\label{sec:limitations}

Chiyoda is currently a stylized research simulator, not an operational
evacuation planner. The station layouts are simplified grids unless imported
from richer geometry. Hazard physics are intentionally lightweight and should
not be read as validated CBRN dispersion models. Physiological impairment uses
piecewise exposure curves that are useful for comparative simulation but not
medical prediction.

The behavioral model is also a hypothesis generator rather than a calibrated
description of a specific crowd. Agents use simplified BDI-style intentions,
local gossip, panic/anxiety states, and familiarity cohorts. These mechanisms
are valuable because they make assumptions inspectable, but the paper must not
claim real-world predictive validity without calibration against evacuation
drills, incident records, controlled VR experiments, or pedestrian trajectory
data.

Finally, entropy is not equivalent to safety. The point of the paper is
precisely that uncertainty reduction can be harmful when it produces synchronized
movement. The proposed information-safety metrics are therefore comparative
diagnostics, not universal welfare functions.

A natural extension is adaptive message generation, including learned or
language-model-mediated guidance. That extension should not be evaluated by
message fluency alone. Under this paper's framing, any generated message is a
safety-control action: it must be replayable, validated against the simulated
state, prevented from inventing unavailable exits or stale hazard claims, and
measured by downstream exposure, queueing, exit balance, and evacuation
outcomes. Without those controls, richer language could make a policy more
persuasive while also amplifying harmful convergence.

\subsection{LLM Extension Protocol}
\label{subsec:llm-extension-protocol}

The repository now includes an optional LLM-guidance policy scaffold, but it is
not part of the completed empirical claims in this paper. Its purpose is to
support a controlled follow-on question: whether generated evacuation messages
add safety value beyond deterministic static, global, entropy-targeted, and
bottleneck-aware baselines. The protocol deliberately treats an LLM as a
bounded message proposer. The simulator exposes a structured state summary,
requires a structured message containing recommended exits, avoided exits,
hazard references, confidence, and optional abstention, and validates the
proposal before it can affect agent beliefs.

This design keeps the extension compatible with the core contribution. A live
API call is not allowed to be the experimental evidence by itself. Paper runs
must use cached or deterministic replay, record the prompt-relevant state and
model metadata, and report rejected messages separately from accepted
interventions. The correct comparison is therefore not whether LLM text sounds
more natural, but whether a validated generated-message policy improves ISE or
reduces HCI under the same information budget and target-selection constraints
as deterministic baselines.

\subsection{Medium LLM Extension Result}
\label{subsec:llm-medium-result}

A first medium-scale LLM extension study applies this protocol to 80 runs:
eight policy variants across ten random seeds. The study includes the
deterministic baselines, a deterministic template generator, three live OpenAI
prompt/validator ablations, and a replay-only verification variant. The live
LLM policies make only eight intervention attempts per run, compared with
32--72 interventions for the deterministic communication policies.

\begin{table}[t]
\centering
\small
\resizebox{\columnwidth}{!}{%
\begin{tabular}{lrrrr}
\toprule
Policy & Evacuated & ISE & HCI & Recipients \\
\midrule
Static beacon & 11.0 & 0.0145 & 6.81 & 503.7 \\
Entropy targeted & 9.2 & 0.0079 & 8.10 & 1064.3 \\
Bottleneck avoidance & 9.4 & 0.0124 & 7.55 & 585.7 \\
Template safety & 9.0 & 0.0155 & 7.86 & 446.6 \\
OpenAI state-only & 9.0 & 0.0414 & 8.39 & 132.3 \\
OpenAI safety-strict & 9.9 & 0.0491 & 7.84 & 133.8 \\
OpenAI entropy-lenient & 10.0 & 0.0429 & 8.76 & 138.8 \\
Replay safety-strict & 9.9 & 0.0491 & 7.84 & 133.8 \\
\bottomrule
\end{tabular}
}
\caption{Medium LLM extension summary over ten seeds. LLM guidance improves
information-safety efficiency mainly by operating under a much smaller
intervention budget, but it does not yet reduce harmful convergence relative
to the strongest conservative baseline.}
\label{tab:llm-medium-extension}
\end{table}

The result is useful but deliberately bounded. The live OpenAI variants
achieve higher information-safety efficiency than the deterministic policies
under this intervention budget, and replay exactly matches the safety-strict
live variant. However, the harmful-convergence index remains high. The
extension therefore supports a cautious paper claim: validated generated
messages can be made replayable and efficient, but richer adaptive language
does not by itself solve crowd synchronization risk. The next experiment
should test whether LLM guidance can reduce HCI under controlled target
selection, hazard severity, and population familiarity regimes.

\section{Conclusion}
\label{sec:conclusion}

Chiyoda reframes evacuation as an information-control problem. In a
hazard-coupled crowd, emergency communication changes beliefs, beliefs change
routes, and route convergence changes exposure and congestion. The framework
therefore measures not only whether information spreads, but whether spreading
it at a particular time, place, credibility, and radius improves safety.

The current empirical package supports that framing, but it also narrows the
claim. Static beacons are the strongest current baseline by
information-safety efficiency, global broadcast shows why reach is not enough,
and entropy-targeted intervention exposes how belief improvement can coexist
with weak aggregate completion. The next step is not to claim operational
readiness, but to test external validity across layouts, hazards, and
population mixes while preserving the safety-control interpretation developed
here.


\bibliographystyle{plain}
\bibliography{references}

\end{document}
